\chapter{Conclusion}
\label{sec:conclusion}

This thesis has investigated whether mechanistic structures described
in Anthropic's \textit{On the Biology of a Large Language Model}
transfer to \texttt{Qwen3-4B-Instruct}, an independently trained
open-access model, and has extended that investigation through
concept localisation and knowledge editing.

The reproducibility study confirmed that lookup
table features analogous to those in Claude 3.5 Haiku are present
in Qwen3-4B, with a comparable positional structure.
Attribution graphs over the linearised transcoder model find
the same high-level circuit topology, and the multilingual experiment
demonstrates that operation, operand, and output-language components
can be independently intervened on, suggesting that the mechanistic
organisation identified by Lindsey et al.\ is not specific to a
single model family.

The concept-localisation study reveals that the geometry of the concept representation
 reflects the algebraic structure of the task, by revealing mathematical 
 structure in the features involved in the internal computation. 
 Approximately dichotomous concepts such as carry and GCD
divisibility produce more localised directions than
residue-class membership, whose contrast does not yield a stable binary
direction, and is therefore harder to intervene on, given contrastive methods.

The soft prompting experiments show that learned prefix embeddings
improve model predictions on the targeted prompts, but the learned
directions are not interpretable and do not correspond to identifiable
concepts in the residual stream.

Finally, these results suggest that mechanistic structure identified
in one model family transfers to independently trained instruction-tuned
models, that concept geometry varies with the algebraic structure of the
task, and that soft interventions can be effective without producing
interpretable representations. Future work should examine whether the
transcoder features identified via delta projection form a sufficient
circuit for the concept, and investigate causal interventions on the
arithmetic computation circuit across standard mathematical benchmarks.
